\chapter{Repair, Cost, and the Economics of Maintaining a Distinction}
\label{ch:repair-economics}

\begin{quotation}
\noindent\textit{Synopsis.} This chapter imports the repair-theoretic
machinery of \emph{Conservation of Ambiguity}: repair cost, ambiguity, and
interference as quantitative descriptions of how finite budgets constrain
distinguishability. The Redistribution Theorem is restated and carefully
scoped as a description of how a single observer reallocates finite
resources across distinctions within one repair economy at one moment. The
chapter's key emphasis, on which the rest of the monograph depends, is that
ambiguity conservation is not that theorem's central result; the deeper
conserved quantity is the repair budget itself, and ambiguity conservation
is recovered only as a restrictive limiting case.
\end{quotation}

\section{Repair Cost as the Cost of Preserving a Minimal Pair}

Chapter~\ref{ch:distinctions} defined a named distinction $d$ via a
minimal pair $\{x_d, x_{\neg d}\}$ (\cref{def:minimal-pair}) and defined
$d$ as discarded by a rendering precisely when that pair collapses
(\cref{def:discard-set}). This chapter's imported apparatus speaks
instead of the \emph{repair cost} of $d$, and it is worth stating plainly,
before using the two vocabularies interchangeably as every later chapter
does, exactly how they connect.

\begin{definition}[Repair Cost, Specialized]
\label{def:repair-cost-specialized}
For a named distinction $d$ with minimal pair $\{x_d, x_{\neg d}\}$
subject to a perturbation process (a transformation, a transmission
channel, a compression step), the repair cost $\repcost{d}$ is the cost of
maintaining $x_d \not\sim x_{\neg d}$ -- that is, of keeping the pair told
apart -- under that perturbation.
\end{definition}

\cref{def:repair-cost-specialized} is a specialization of
Chapter~\ref{ch:repair-economics}'s imported notion of repair cost to the
minimal-pair vocabulary of Chapter~\ref{ch:distinctions}, and it is what
licenses Chapter~\ref{ch:budget-relay}'s otherwise unexplained move of
treating $\repcost{d}$ and $\discard{r}$ as talking about the same
underlying phenomenon: a distinction is discarded by a rendering exactly
when insufficient repair investment was made, relative to that
rendering's perturbation, to keep its minimal pair apart.

\section{Repair Cost, Sharpness, and Ambiguity}

\imported{Conservation of Ambiguity}

\begin{definition}[Distinction Sharpness]
\label{def:sharpness}
For $d \in \restdomain$ a repair-stable distinction, the sharpness of $d$
is $\sigma(d) = \repcost{d}$, the repair cost required to maintain $d$
under admissible perturbation.
\end{definition}

\begin{definition}[Operational Ambiguity]
\label{def:ambiguity}
The ambiguity of a distinction is a monotone decreasing function of its
repair cost; the simplest choice is $\amb{d} = 1/\repcost{d}$.
\end{definition}

\begin{remark}[Ambiguity's Functional Form Is a Convention, Not a Law]
\label{rem:ambiguity-convention}
\emph{Conservation of Ambiguity} states this explicitly, and it is worth
repeating rather than letting the simplest choice harden into an
assumption: any strictly decreasing monotone transformation of
$\repcost{d}$ produces an equivalent \emph{ordering} of distinctions by
ambiguity, so $\amb{d} = 1/\repcost{d}$ is adopted here for concreteness
and ease of computation in worked examples, not because it is uniquely
forced by anything proved. Nothing in this monograph depends on the
specific functional form; every result stated in terms of $\amb{d}$ could
be restated, without loss, in terms of the ordering it induces.
\end{remark}

Ambiguity, under \cref{def:ambiguity}, is not ignorance, incomplete
information, or subjective uncertainty. It measures the degree to which a
distinction has been left under-maintained: a distinction is ambiguous
because insufficient repair capacity has been committed to stabilizing it
against perturbation, not because information about it is unknowable in
principle.

\section{Repair Interference}

\begin{definition}[Repair Interference]
\label{def:interference}
For $d_i, d_j \in \restdomain$, the repair interference
$\interf{d_i}{d_j}$ is the change in repair cost of $d_j$ induced by an
infinitesimal increase in repair investment devoted to $d_i$. Positive
interference indicates repair competition; negative interference
indicates repair synergy; zero interference indicates operational
independence.
\end{definition}

\section{The Redistribution Theorem}

\begin{theorem}[Redistribution Theorem]
\label{thm:redistribution}
\imported{Conservation of Ambiguity}
Let $\Phi = \amb{\cdot} + \lambda \repcost{\cdot} + \mu I(\cdot,\cdot)$ be
the repair budget functional of a finite observer, with total budget
$\Phi_{\mathrm{tot}}$ held fixed. Then every admissible variation
satisfying $\delta \repcost{d_1} > 0$ for some distinction $d_1$ induces
compensating variations throughout the remaining distinction space such
that $\delta\Phi = 0$. Consequently, no local increase in distinction
sharpness can occur without global redistribution of ambiguity, repair
cost, or repair interference.
\end{theorem}

\begin{proofsketch}
\imported{Conservation of Ambiguity} The result follows from constrained
variation of the repair budget functional $\Phi$ under fixed total
budget. Since the observer possesses no additional maintenance capacity,
every increase in repair investment devoted to one region of distinction
space must appear elsewhere as a corresponding change in ambiguity,
repair cost, or repair interference. The full proof, including the three
interference regimes (competition, synergy, independence), is given in
\emph{Conservation of Ambiguity} and is not repeated here.
\end{proofsketch}

\textbf{Scope note.} \cref{thm:redistribution} governs a single observer's
competing distinctions at a single moment: a \emph{synchronic} result. It
says nothing, on its own, about what happens to a single distinction's
status as it is carried across a sequence of independently-budgeted
stages over time. That extension is the subject of
Chapter~\ref{ch:synchronic-to-diachronic} and Part~\ref{part:budget-relay}.

\begin{corollary}[Ambiguity Conservation as a Limiting Case]
\label{cor:ambiguity-limiting-case}
\imported{Conservation of Ambiguity}
If repair costs and interference remain fixed ($\delta\repcost{\cdot} =
\delta I(\cdot,\cdot) = 0$), then $\delta \amb{\cdot} = 0$: ambiguity
conservation is recovered as a special case of repair-budget conservation.
In realistic systems, sharpening one distinction almost always changes
both repair costs and interference relationships, so ambiguity alone is
generally \emph{not} conserved; the conserved quantity is the complete
repair budget $\Phi$.
\end{corollary}

This corollary is the reason this monograph does not treat itself as a
corollary of \cref{thm:redistribution}: the earlier framing of the
research program that produced this book -- ``ambiguity is conserved,
merely relocated'' -- is precisely the special case
\cref{cor:ambiguity-limiting-case} shows to be exceptionally restrictive
and generally false. The new work required here
(Part~\ref{part:budget-relay}) does not inherit its central claim from
this theorem; it must be argued independently.

\section{A Worked Instance, Before It Becomes a Relay}
\label{sec:redistribution-worked}

Chapter~\ref{ch:budget-relay} introduces a running numerical example: five
distinctions $d_1, \ldots, d_5$ with repair costs $\repcost{d_1}=1$,
$\repcost{d_2}=2$, $\repcost{d_3}=3$, $\repcost{d_4}=2$,
$\repcost{d_5}=4$. There, the costs are held fixed and a \emph{budget}
$\beta_i$ varies from stage to stage of a relay. Here, before any relay
exists, it is worth seeing what \cref{thm:redistribution} says about the
same five distinctions considered as the simultaneous concern of a
\emph{single} observer with one fixed total budget
$\Phi_{\mathrm{tot}}$ -- the case the theorem actually governs.

Suppose the observer initially funds all five distinctions exactly at
their stated costs, for a total commitment of $1+2+3+2+4 = 12$, and
suppose $\Phi_{\mathrm{tot}} = 12$ exactly, leaving no slack. Now suppose
the observer wishes to sharpen $d_3$ -- to invest further in it beyond its
baseline cost, driving $\repcost{d_3}$ up and $\amb{d_3}$ down, perhaps
because $d_3$ has become newly important. \cref{thm:redistribution}
states that this cannot happen for free: with $\Phi_{\mathrm{tot}}$ fixed
at $12$ and already fully committed, $\delta \repcost{d_3} > 0$ forces a
compensating $\delta \Phi = 0$ elsewhere in the functional. If
interference is negligible ($I \approx 0$, as in the worked relay of
\cref{sec:worked-relay}), the compensation must come directly out of the
repair cost or ambiguity terms of the remaining four distinctions --
concretely, some combination of $d_1, d_2, d_4, d_5$ must have its own
funding reduced by exactly the amount $d_3$'s sharpening consumed, raising
their ambiguity in turn.

This is the single-observer, single-moment phenomenon
\cref{thm:redistribution} describes, and it is worth contrasting directly
with what happens to the \emph{same five numbers} once
Chapter~\ref{ch:budget-relay} turns them into a two-stage relay with
$\beta_0 = 4$ and $\beta_1 = 3$. There, $d_3$'s low priority at stage $0$
does not trigger any compensating adjustment to $d_1$ or $d_2$'s funding
elsewhere in a shared $\Phi_{\mathrm{tot}}$, because no such shared
quantity exists across stages (\cref{def:budget-relay}): $d_3$, $d_4$, and
$d_5$ are simply dropped, permanently, with nothing anywhere in the
system compensating for their loss. The redistribution
\cref{thm:redistribution} guarantees within one economy has no
counterpart at all across the boundary between $\beta_0$ and $\beta_1$.
Chapter~\ref{ch:synchronic-to-diachronic} argues for this contrast in
prose; \cref{prop:no-global-conserved-quantity} there states it as a
precise claim about why no functional analogous to $\Phi$ can be defined
across a relay's stages in the first place.
